\documentclass[10pt,a4paper]{amsart}
\usepackage[margin=19mm]{geometry}
\usepackage{amsmath,amssymb,mathtools}
\usepackage[scr=rsfs,cal=euler]{mathalfa}
\usepackage{xfrac}
\usepackage[unicode,hidelinks,bookmarks=false]{hyperref}
\newcommand{\ie}{i.e.\ }
\newcommand{\cf}{cf.\ }
\newcommand{\resp}{respectively\ }
\newcommand{\Subsection}{\subsection}
\providecommand{\nobreakdash}{\nobreak\hbox{-}}
\setlength{\emergencystretch}{2em}
\multlinegap0pt
\usepackage{amssymb}
\usepackage{mathtools}
\usepackage{xcolor}
\newtheorem{lemma}{Lemma}
\newcommand{\UU}{\mathcal U}
\DeclareMathOperator{\dist}{dist}
\newcommand{\mm}{\mathfrak m}
\title{A one-dimensional coherent domain\\ with non-Pr\"ufer integral closure}
\author{Viet-Hoang Tran, Thieu N. Vo, Tan M. Nguyen}
\date{Source: arXiv:2609.21013v1 (CC BY 4.0)}
\begin{document}
\maketitle

\noindent\emph{Source and licence.} A one-dimensional coherent domain\\ with non-Pr\"ufer integral closure. Version arXiv:2609.21013v1 (CC BY 4.0), distributed by arXiv under the Creative Commons Attribution 4.0 International licence, http://creativecommons.org/licenses/by/4.0/, as recorded in the arXiv licence metadata for this exact version and shown on the abstract page https://arxiv.org/abs/2609.21013v1. Full licence notice: \texttt{licenses/arxiv-2609.21013-CC-BY-4.0.txt}.\par\smallskip

\noindent\textit{Editorial scope.} Counted unit: the article's lemma lem:separation (source0.tex lines 327--333). The article's complete original proof of it is delivered with it (source0.tex lines 334--385, one \texttt{proof} environment of 52 lines). No mathematics is written, repaired or re-derived: the statement and the proof are the article's own lines, in the article's own notation, and no retained line is substituted or re-flowed.  No further source-local context is needed: every cross-reference of every retained line resolves inside the two delivered spans.  Nothing was omitted from any retained span, so the package carries no omission record.  No retained line contains a citation, so the wrapper carries no bibliography and no bibliography entry.  No retained line contains a figure environment, an image-inclusion command or drawing code, so no image is delivered and the package declares no asset dependency.  Editorial formatting only, in addition: an AMS \texttt{amsart} preamble that restates the article's own declarations and macros used by the retained lines, namely the environments \texttt{lemma} on separate counters, together with the standard packages those lines need.  The retained environments print in this document as Lemma 1 in order of appearance, whereas the article prints the same items with the numbering of its own full document.  The complete native source of the article as distributed in the arXiv e-print for this exact version is bundled unchanged as \texttt{sources/210924/source0.tex}.
\begin{lemma}\label{lem:separation}
Every polynomial in $R[T]$ vanishing at $z$ has all coefficients in
$\mm$. Consequently
\[
 R[z]/\mm R[z]\cong(R/\mm)[T],\qquad \dim R[z]\geq2.
\]
\end{lemma}

\begin{proof}
Suppose $\sum_{i=0}^d c_i z^i=0$ and some $c_i\notin\mm$;
necessarily $d\geq1$. Clear a common denominator to write
$c_i=\bar b_i/\bar s$, with $b_i\in A$ and $s\notin Q$.
Then $J=\{i:b_i\notin Q\}\ne\varnothing$ and, since
$A/P\hookrightarrow F$,
\begin{equation}\label{eq:error}
 h=\sum_{i=0}^d b_i v^i u^{d-i}\in P.
\end{equation}
For this fixed degree, irrationality gives
\[
 \delta=\min_{1\leq k\leq d}\dist(k\alpha,\mathbb Z)>0.
\]
Choose a common $\epsilon>0$ with
$\nu_n(b_{i,n})>\epsilon n^2$ almost everywhere for all $i\notin J$
(take $\epsilon=1$ if there are none), and fix
$0<\eta<\min\{\epsilon/4,\delta/4,1/4\}$.
On a common set $E\in\UU$ these lower bounds hold,
$\nu_n(b_{i,n})\leq\eta n^2$ for every $i\in J$, and
$d/n<\min\{\delta/4,\eta\}$.

For $n\in E$ and $i\in J$, write the leading exponent as
$e_{i,n}=q_{i,n}n+r_{i,n}$, $0\leq r_{i,n}<n$.
The semigroup criterion gives
\[
 0\leq r_{i,n}/n\leq q_{i,n}/n\leq e_{i,n}/n^2\leq\eta.
\]
The shifted orders $e_{i,n}+id_n$ are distinct for $i\in J$.
Otherwise equality for $i\ne j$, using
$q_{i,n}-q_{j,n}\in\mathbb Z$, would give
\[
 \delta\leq\dist((i-j)\alpha,\mathbb Z)
 \leq\frac{|r_{i,n}-r_{j,n}|}{n}
      +|i-j|\left|\alpha-\frac{d_n}{n}\right|
 <\eta+\frac dn<\frac\delta2.
\]
For $i\in J$, the order of $b_{i,n}t^{id_n}$ is less than
$2\eta n^2<\epsilon n^2$; for $i\notin J$ it exceeds
$\epsilon n^2$. Thus $g_n=\sum_i b_{i,n}t^{id_n}$ has a unique
summand of least order, and $\nu_n(g_n)<2\eta n^2<n^2$.
Since $h_n=t^{dn^2}g_n$, this gives
$\nu_n(h_n)<(d+1)n^2$ on $E$, contradicting \eqref{eq:error}.

The evaluation kernel lies in $\mm R[T]$, giving the quotient
isomorphism. The primes $(0)\subsetneq(T)$ of $(R/\mm)[T]$ yield
the prime chain
\[
 0\subsetneq\mm R[z]\subsetneq\mm R[z]+zR[z].
\]
The first inclusion is strict since $0\ne\bar u\in\mm R[z]$;
hence $\dim R[z]\geq2$.
\end{proof}

\end{document}
